\begin{problem}{TST 2026 - Bài 4}{standard input}{standard output}{2.5 seconds}{1024 megabytes}

Alice có một dãy số $A$ gồm $N$ phần tử, các phần tử được đánh số từ $0$ đến $N-1$. Ban đầu, với mọi $i$ từ $0$ đến $N-1$, $A_i=i$. Ngoài ra, Alice có một danh sách gồm $M$ thao tác, được đánh số từ $0$ đến $M-1$. Với thao tác thứ $i$ ($0\le i\le M-1$):
\begin{itemize}
   \item Alice sẽ chọn phần tử \textbf{có giá trị $H_i$} và chuyển nó về đầu dãy số. Các phần tử đứng trước $H_i$ sẽ lùi sang bên phải một vị trí.
\end{itemize}

Alice đưa hai dãy số $A$ và $H$ cho Bob, sau đó hỏi Bob $Q$ câu hỏi. Với câu hỏi thứ $i+1$ ($0\le i\le Q-1$):
\begin{itemize}
   \item Alice đưa cho Bob ba số nguyên $L_i,R_i,P_i$.
   \item Alice sẽ thực hiện tất cả các thao tác theo thứ tự, nhưng không thực hiện các thao tác nằm trong đoạn $[L_i;R_i]$ \textbf{trên dãy ban đầu}.
   \item Bob cần xác định chính xác vị trí của phần tử có giá trị $P_i$ sau khi Alice đã thực hiện $M-(R_i-L_i+1)$ thao tác trên dãy ban đầu.
\end{itemize}

\textbf{Yêu cầu:} Cho hai dãy số $A,H$ và $Q$ câu hỏi, hãy giúp Bob viết một chương trình trả lời các câu hỏi do Alice đưa ra.

\InputFile
\textbf{Chi tiết cài đặt}

Thí sinh cần cài đặt hàm sau:

\begin{lstlisting}
vector<int> solve(int N, int M, int Q, const vector<int>& H, const vector<int>& L,
                                       const vector<int>& R, const vector<int>& P);
\end{lstlisting}

\begin{itemize}
    \item $N$: Số lượng phần tử trong dãy số $A$.
    \item $M$: Số thao tác trong dãy thao tác $H$ của Alice.
    \item $Q$: Số câu hỏi Alice hỏi Bob.
    \item $H$: Mảng gồm $M$ phần tử, trong đó $H_i$ ($0\le i\le M-1$) đại diện cho giá trị của phần tử mà Alice sẽ dịch chuyển trong thao tác thứ $i$.
    \item $L, R$: Hai mảng gồm $Q$ phần tử, trong đó $L_i$ và $R_i$ ($0\le i\le Q-1$) đại diện cho đoạn truy vấn mà Alice sẽ không thực hiện khi xét câu hỏi thứ $i+1$.
    \item $P$: Mảng gồm $Q$ phần tử, trong đó $P_i$ ($0\le i\le Q-1$) đại diện cho giá trị của phần tử Bob cần xác định vị trí trong câu hỏi thứ $i+1$.
    \item Hàm này cần trả về một mảng $ans$ bao gồm $Q$ phần tử, trong đó $ans_i$ ($0\le i\le Q-1$) là đáp án của câu hỏi thứ $i+1$.
    \item Hàm này sẽ được gọi \textbf{một lần duy nhất} với mỗi test của bài.
\end{itemize}

\textbf{Lưu ý:}
\begin{itemize}
    \item Chương trình của bạn cần phải khai báo thư viện \texttt{sgamelib.h} ở đầu.
    \item Trong chương trình, thí sinh được phép cài đặt các biến, các hàm toàn cục, nhưng không được phép có hàm \texttt{main()}.
   \item Chương trình của bạn không được tương tác trực tiếp với đầu vào chuẩn (stdin) và đầu ra chuẩn (stdout).
\end{itemize}

\textbf{Ràng buộc}

\begin{itemize}
    \item $1\le N, M, Q \le 2\times 10^5$
    \item $0\le H_i \le N-1$ với mọi $i$ từ $0$ đến $M-1$
    \item $0\le L_i \le R_i \le M-1$ với mọi $i$ từ $0$ đến $Q-1$
    \item $0\le P_i \le N-1$ với mọi $i$ từ $0$ đến $Q-1$
\end{itemize}

\Scoring
\begin{center}
  \begin{tabular}{|c|c|l|} \hline
    \textbf{Subtask} & \textbf{Điểm} & \textbf{Giới hạn} \\ \hline
    1 & $5$  & $1\le N, M, Q \le 500$ \\ \hline
    2 & $10$ & $1\le N, M, Q \le 5000$ \\ \hline
    3 & $33$ & $\sum_{i=0}^{Q-1} (R_i-L_i) \le 2\times 10^5$ \\ \hline
    4 & $52$ & Không có ràng buộc gì thêm \\ \hline
  \end{tabular}
\end{center}


\Example

\begin{example}
\exmpfile{example.01}{example.01.a}%
\end{example}

\Note
\textbf{Example}

Xét lần gọi hàm sau:

\begin{lstlisting}
solve(4, 3, 2, \{2, 0, 3\}, \{1, 0\}, \{1, 0\}, \{1, 1\})
\end{lstlisting}

Xét câu hỏi đầu tiên, $L_i=1, R_i=1, P_i=1$. Quá trình thực hiện thao tác của Alice sẽ được mô tả như sau:
\begin{itemize}
    \item Ban đầu, $A=\{0,1,2,3\}$.
    \item Alice thực hiện thao tác số $0$, sau thao tác $A=\{2,0,1,3\}$.
    \item Alice không thực hiện thao tác số $1$ vì $1\in [1;1]$.
    \item Alice thực hiện thao tác số $2$, sau thao tác $A=\{3,2,0,1\}$.
\end{itemize}

Vì vị trí của phần tử có giá trị $P_i=1$ là $3$ nên đáp án cho câu hỏi đầu tiên là $3$.

Xét câu hỏi thứ hai, $L_i=0, R_i=0, P_i=1$. Quá trình thực hiện thao tác của Alice sẽ được mô tả như sau:
\begin{itemize}
    \item Ban đầu, $A=\{0,1,2,3\}$.
    \item Alice không thực hiện thao tác số $0$ vì $0\in [0;0]$.
    \item Alice thực hiện thao tác số $1$, sau thao tác $A=\{0,1,2,3\}$ (phần tử giá trị $0$ chuyển lên đầu).
    \item Alice thực hiện thao tác số $2$, sau thao tác $A=\{3,0,1,2\}$.
\end{itemize}

Vì vị trí của phần tử có giá trị $P_i=1$ là $2$ nên đáp án cho câu hỏi thứ hai là $2$.

Như vậy, hàm trên cần trả về kết quả là $\{3,2\}$.

\end{problem}

